\section{Conclusion}
\label{sec:conclusion}
Following the project from the exact model to held-out data, four lessons stood out.
\begin{itemize}
  \item \emph{Accuracy is the wrong target when errors cost different amounts.} The most accurate rule
        was the most expensive one in the exact model (Section~\ref{sec:m6}), in the cost sweep
        (Section~\ref{sec:m2}) and on held-out data (Section~\ref{sec:m4}). With $\cfp=10$ the right cut
        is $10/11$, not $0.5$.
  \item \emph{The rule is only as good as the probabilities.} Two models with the same accuracy can
        have very different costs (Section~\ref{sec:m5}). Logistic regression gave the best-calibrated
        probabilities and the lowest costs on Spambase; the Gaussian models needed $\log(1+x)$ features,
        shrinkage and re-calibration to come close (Sections~\ref{sec:m3} and~\ref{sec:m5}).
  \item \emph{A change in the spam rate is cheap to fix.} One shift of the log-odds more than halved
        the cost when spam fell from $39\%$ to $5\%$, even with the new rate estimated from unlabelled
        emails (Section~\ref{sec:m1}).
  \item \emph{Inspection is worth paying for when it is cheap.} At $\ci=0.3$ it cut the cost on
        Spambase by $0.10$--$0.12$ per email, inspecting a quarter of the emails. It is never used once
        $\ci\ge10/11$ (Section~\ref{sec:m4}).
\end{itemize}
\emph{Limitations.} The inspector is assumed to be always right, the costs are our own choice, and
each Spambase test set has only about $920$ emails, so differences within the standard errors should
not be read as real.
\code{python -m project4.run\_all} regenerates every number in this report.
